\FloatBarrier
\section{How Voltages Set and Preserve Model Parameters}
\label{app:voltage-response}
\label{sec:device-parameter-dependence}

A gate electrode changes the potential over a finite part of the
device. Its voltage can therefore shift a dot level, reshape a
localized orbital and change tunneling through a nearby barrier
at the same time. Moving an addition threshold in
Section~\ref{sec:electron-number-control} requires knowing its voltage
response. Tuning the exchange barrier at fixed thresholds in
Section~\ref{sec:exchange-virtual-tunneling} requires coordinating
several electrodes. This appendix derives the parameter responses
and then solves for the voltage combinations that produce a specified
change. The same construction is extended to an array, where the
selection problem of Section~\ref{sec:platform-array-selection}
requires constraints on affected neighbours as well as on the target.

\subsection{How much does a gate voltage move an energy level?}

Consider a concrete operation: increase the voltage
on the metal electrode above a dot. The resulting electric
potential change lowers the electron potential energy.
The oxide, surrounding conductors, and electron position determine
how strongly the dot responds.

Let $\mathbf V$ denote the collection of gate voltages.
For a small change around a specified operating point, define
the dimensionless response $\eta_g(\mathbf r)$ by
\begin{equation}
 \delta\Phi_0(\mathbf r)
    =\sum_g\eta_g(\mathbf r)\,\delta V_g,
 \qquad
 \eta_g(\mathbf r)=
       \frac{\partial\Phi_0(\mathbf r)}{\partial V_g}.
 \label{eq:device-voltage-response}
\end{equation}
It is obtained by changing a gate boundary voltage in the
electrostatic problem. It can be calculated by differencing
nearby solutions; if the charge response is nonlinear, it
depends on the chosen operating point.

For a normalized, nondegenerate eigenstate $w_a$ of an isolated
dot, assume initially that voltage dependence enters only through
$-e\Phi_0$. Differentiating its eigenvalue equation gives the
Hellmann--Feynman relation
\begin{equation}
 \frac{\partial\epsilon_a}{\partial V_g}
    =\bra{w_a}\frac{\partial h_1}{\partial V_g}\ket{w_a}
    =-e\int d^3r\,|w_a(\mathbf r)|^2\eta_g(\mathbf r).
 \label{eq:device-hellmann-feynman}
\end{equation}
Terms containing derivatives of $w_a$ cancel by normalization
and the eigenvalue equation. Thus the electron samples the
gate's potential response over its probability distribution.

The resulting voltage-to-energy conversion is called a
\emph{lever arm}. Using a dimensionless convention,
\begin{equation}
 \alpha_{ag}\equiv-\frac1e
           \frac{\partial\epsilon_a}{\partial V_g},
 \qquad
 \delta\epsilon_a=-e\sum_g\alpha_{ag}\delta V_g.
 \label{eq:device-lever-arm}
\end{equation}
Some literature instead quotes $e\alpha_{ag}$, an energy per
voltage. The convention must be stated.
If voltage also changes interface couplings or other terms in
$h_1$, their derivatives must be included in
Eq.~\eqref{eq:device-hellmann-feynman}.

An addition threshold depends on the energies of two many-electron
states. Its voltage response is
\begin{equation}
 \frac{\partial\mu_N}{\partial V_g}
    =\frac{\partial E_N}{\partial V_g}
      -\frac{\partial E_{N-1}}{\partial V_g}.
 \label{eq:device-addition-response}
\end{equation}
For example, a second electron in a frozen single orbital has
addition energy $\epsilon+U$, so its shift contains both
$\partial\epsilon/\partial V_g$ and $\partial U/\partial V_g$.
This is why converting a charge transition into a model
parameter requires identifying which energy difference is
being measured.

\subsection{Changing tunneling and repulsion}

Changing the electrode between two dots changes the barrier
through which their wave functions extend. Its most direct
effect is therefore on $t_{12}$.

A one-dimensional barrier gives a useful estimate. Let $s$
measure distance along a tunneling path, $V_b(s)$ be its potential
energy, and $E_0$ the approximate bound-state energy. In the
forbidden interval $s_1<s<s_2$, where $V_b>E_0$, an effective
mass $m_b$ gives the leading exponential dependence
\begin{equation}
 |t_{12}|\ \propto\
 \exp\!\left[-\int_{s_1}^{s_2}
       \frac{\sqrt{2m_b[V_b(s)-E_0]}}{\hbar}\,ds\right].
 \label{eq:device-wkb-hopping}
\end{equation}
This WKB estimate explains why a modest change in barrier height
or width can strongly change tunneling. The prefactor and the
actual multidimensional coupling should be calculated from the
wave functions. Here $t_{12}$ is a coherent Hamiltonian matrix
element with units of energy.

Gate voltages also change the orbital sizes and positions,
so they change the interaction integrals.
For a family of similarly shaped orbitals of radius $\ell$
in a uniform dielectric,
$U$ scales roughly as $e^2/(4\pi\epsilon_{\mathrm d}\ell)$.
Spreading an orbital out reduces the cost of placing two
electrons in it. For well-separated dots at distance $R$,
the unscreened direct interaction scales as
$V_{12}\sim e^2/(4\pi\epsilon_{\mathrm d}R)$.
The scaling estimates apply to unscreened interactions in a
uniform dielectric. Nearby metals modify them through the
screening described in Eq.~\eqref{eq:device-image-screening}.

Consequently, geometry affects more than the shape of
$V_{\mathrm{conf}}$. Oxide thickness affects gate control and
screening; dot separation affects tunneling and interdot
repulsion; interface structure affects the valley states.
The integrals must be recalculated, or their changes justified
as negligible, when these inputs change appreciably.

\subsection{Why one physical control changes several model parameters}

Consider a two-dot device with a plunger above each dot and
an electrode between them. Changing the middle electrode may
increase tunneling while also shifting both dot energies.
This simultaneous response to one gate voltage is called
\emph{gate cross-coupling}.
To keep the energies fixed, one can simultaneously adjust the
two plungers. Near a reference voltage setting, the required
compensation obeys
\begin{equation}
 \delta\epsilon_i
    =\sum_g
       \frac{\partial\epsilon_i}{\partial V_g}\delta V_g=0,
 \qquad i=1,2.
 \label{eq:device-compensation}
\end{equation}
For the Hubbard on-site energies, evaluate the derivatives
including the voltage dependence of the chosen localized basis.
Independent plunger responses allow these two linear equations
to determine compensating voltages locally. The interaction
parameters $U_i,V_{12}$ and orbital shapes can still vary along
this compensated voltage change.

The exchange coupling inherits these simultaneous changes.
For two dots with equal $U$, define
$D=U-V_{12}$ and $\varepsilon=\epsilon_1-\epsilon_2$.
Writing Eq.~\eqref{eq:exchange-detuning} with
$t_c=|t_{12}|$ and $V=V_{12}$ gives
\begin{equation}
 J_{\mathrm{se}}(\varepsilon)
    =\frac{4|t_{12}|^2D}{D^2-\varepsilon^2},
 \qquad D>|\varepsilon|,\quad
 |t_{12}|\ll D-|\varepsilon|.
 \label{eq:device-exchange-control}
\end{equation}
Holding $t_{12}$ and $D$ fixed, the first derivative with respect
to detuning vanishes at $\varepsilon=0$.
This explains why operation near equal dot energies can reduce
first-order sensitivity to detuning fluctuations~\cite{Reed2016}.
The sensitivity to changes in the tunnel amplitude and
interaction energy follows, at this point, from the first-order relation
\begin{equation}
 \frac{\delta J_{\mathrm{se}}}{J_{\mathrm{se}}}
       =2\frac{\delta|t_{12}|}{|t_{12}|}
          -\frac{\delta D}{D}.
 \label{eq:device-exchange-sensitivity}
\end{equation}
A voltage that changes the barrier can therefore tune exchange
even when the dot energies remain aligned.

The spin splitting can change as well. If a dot moves through
the field of a micromagnet, its mean field changes. For a scalar,
position-independent electron $g$ factor $g_{\mathrm e}$ and a
retained normalized orbital $w_i$, the Zeeman vector is
\begin{equation}
 \mathbf b_i
    =g_{\mathrm e}\mu_B\int d^3r\,|w_i(\mathbf r)|^2\mathbf B(\mathbf r).
 \label{eq:device-averaged-zeeman}
\end{equation}
Here $\mu_B$ is the Bohr magneton and $\mathbf B$ is the applied
plus micromagnet field.
Changing $w_i$ changes $\mathbf b_i$ even when the magnet is
unchanged. This connects the device calculation back to
electrically driven spin control in Section~\ref{sec:transverse-spin-control}.

% Included by appendix-voltage-response.tex.
\subsection{Choosing voltages for specified addition-energy changes}
\label{app:virtual-gates}
\label{sec:operating-points-virtual-gates}

Loading requires moving an addition energy below a reservoir chemical
potential; a compensated exchange pulse instead holds the addition
thresholds fixed while changing tunneling. Because one electrode can
produce both effects, we use the combined response to all varied
voltages to select the desired change. For this local calculation,
the response can be measured without recalculating the full device
potential for every pulse.

The many-electron Hamiltonian identifies what this measured
response represents.
For a small voltage change, write
\[
 \delta H=\sum_j \widehat W_j\,\delta V_j,\qquad
 \widehat W_j=\left.\frac{\partial H}{\partial V_j}\right|_{\mathbf V^0}.
\]
For an isolated, nondegenerate state on a specified charge branch,
the Hellmann--Feynman relation gives
\[
 \frac{\partial(E_N-E_{N-1})}{\partial V_j}
 =\langle\widehat W_j\rangle_N-\langle\widehat W_j\rangle_{N-1}.
\]
This is the addition-energy derivative in
Eq.~\eqref{eq:device-addition-response}, expressed through the
voltage-dependent Hamiltonian.

Let $\mu_i(\mathbf V)$ be the addition energy of dot $i$ for a
specified charge transition. Near reference voltages $\mathbf V^0$,
the response to the $N_g$ varied gate electrodes is
\begin{equation}
 d\mu_i=\sum_{j=1}^{N_g}G_{ij}\,dV_j,\qquad
 G_{ij}=\left.\frac{\partial\mu_i}{\partial V_j}\right|_{\mathbf V^0},
 \qquad d\boldsymbol\mu=G\,d\mathbf V.
 \label{eq:calibration-response-matrix}
\end{equation}
The index $j$ includes plunger and barrier electrodes.
The reservoir contact voltages and magnetic field are held fixed
in the examples below. Together with the chosen charge
configuration, these reference settings define the
\emph{operating point}. Each matrix entry has units of energy
per voltage. For a finite small step, the differential relation
gives $\Delta\boldsymbol\mu\simeq G\Delta\mathbf V$.

Suppose we want to increase $\mu_1-\mu_2$ without changing
$(\mu_1+\mu_2)/2$. The required energy changes have opposite signs,
whereas a pulse on one plunger may shift both energies in the
same direction. We use the measured response matrix to find
voltages whose combined effects give the desired pair of changes.

Consider two dots controlled by two plunger voltages $V_1,V_2$
and one barrier voltage $V_{\mathrm B}$. The full relation is
\begin{equation}
 \begin{pmatrix}d\mu_1\\d\mu_2\end{pmatrix}
 =
 \underbrace{\begin{pmatrix}
 G_{11}&G_{12}&G_{1\mathrm B}\\
 G_{21}&G_{22}&G_{2\mathrm B}
 \end{pmatrix}}_{G}
 \begin{pmatrix}dV_1\\dV_2\\dV_{\mathrm B}\end{pmatrix}.
 \label{eq:calibration-full-gate-response}
\end{equation}

\subsubsection{Obtaining the response matrix}
The control calculation needs all six entries, including those of
the barrier electrode. The charge-transition slopes and
finite-bias measurements provide these entries;
we now assemble them into $G$. On a dot--reservoir transition,
$\mu_i=\mu_{\mathrm R}$. Vary electrode $j$ and adjust the
corresponding plunger $V_i$ to stay on that transition, with the
remaining electrode fixed. The measured slope is
\[
 c_{ij}=\left.\frac{dV_i}{dV_j}\right|_{\mu_i=\mu_{\mathrm R}},
 \qquad j\ne i.
\]
Since $0=G_{ii}dV_i+G_{ij}dV_j$, it gives
$G_{ij}=-G_{ii}c_{ij}$. The four slopes
$c_{12},c_{1\mathrm B},c_{21},c_{2\mathrm B}$ determine the
cross-responses relative to the two plunger responses.
Measure the latter,
$g_1=G_{11}$ and $g_2=G_{22}$, using the finite-bias boundaries
of Section~\ref{sec:measurement-bias-spectroscopy}.
The complete matrix is then
\begin{equation}
 G=\begin{pmatrix}
 g_1&-g_1c_{12}&-g_1c_{1\mathrm B}\\
 -g_2c_{21}&g_2&-g_2c_{2\mathrm B}
 \end{pmatrix}.
 \label{eq:calibration-matrix-from-lines}
\end{equation}
The transition slopes give ratios; the finite-bias measurement
supplies the energy scale. Use nearby smooth segments for the
specified charge sectors, since a change in a neighbouring
occupation can shift a transition discontinuously.

\subsubsection{Solving for the voltage change}
With $G$ known, the requested change $d\boldsymbol\mu^\star$ is the
input and $d\mathbf V$ is the unknown. The control problem is
$G\,d\mathbf V=d\boldsymbol\mu^\star$.
A rank-two $2\times3$ matrix leaves one free voltage direction:
two energy conditions do not uniquely determine three voltages.
To write the solutions explicitly, partition the already measured
matrix as
\[
 G=\begin{pmatrix}G_{\mathrm P}&\mathbf g_{\mathrm B}\end{pmatrix},
 \qquad
 d\mathbf V_{\mathrm P}=\begin{pmatrix}dV_1\\dV_2\end{pmatrix},
\]
where $G_{\mathrm P}$ contains the first two columns and
$\mathbf g_{\mathrm B}$ is the third. If $G_{\mathrm P}$ is
invertible, specifying $dV_{\mathrm B}$ gives
\begin{equation}
 d\mathbf V_{\mathrm P}
 =G_{\mathrm P}^{-1}
       \left(d\boldsymbol\mu^\star-\mathbf g_{\mathrm B}dV_{\mathrm B}\right).
 \label{eq:calibration-virtual-gates}
\end{equation}
Thus keeping the barrier voltage fixed is one possible choice;
the full response matrix also permits solutions that change it.
Coordinated voltage directions chosen to control specified
physical quantities are called \emph{virtual gates}.

\subsubsection{An opposite shift of two addition energies}
For a numerical illustration, take
\[
 G=\begin{pmatrix}-100&-25&-10\\-20&-80&-15\end{pmatrix}
       \frac{\mu\mathrm{eV}}{\mathrm{mV}}.
\]
These are illustrative response values, not experimental data.
Let $\varepsilon_\mu=\mu_1-\mu_2$ and
$\bar\mu=(\mu_1+\mu_2)/2$.
Increasing $\varepsilon_\mu$ by $80\,\mu\mathrm{eV}$ at fixed
$\bar\mu$ requires
$(\Delta\mu_1,\Delta\mu_2)=(40,-40)\,\mu\mathrm{eV}$.
Choosing $\Delta V_{\mathrm B}=0$ gives
\begin{equation}
 \begin{pmatrix}\Delta V_1\\\Delta V_2\\\Delta V_{\mathrm B}\end{pmatrix}
 \simeq
 \begin{pmatrix}-0.56\\0.64\\0\end{pmatrix}\mathrm{mV}.
 \label{eq:calibration-two-gate-example}
\end{equation}
Substitution gives $56-16=40$ and $11.2-51.2=-40$,
in $\mu\mathrm{eV}$, for the two energy changes.
More generally, with $dV_{\mathrm B}=0$,
\begin{equation}
 d\mathbf V_{\mathrm P}
 =G_{\mathrm P}^{-1}
   \begin{pmatrix}1&1/2\\1&-1/2\end{pmatrix}
   \begin{pmatrix}d\bar\mu\\d\varepsilon_\mu\end{pmatrix}.
 \label{eq:calibration-common-differential}
\end{equation}
The difference $\varepsilon_\mu$ concerns the chosen addition
transitions. Its relation to the orbital detuning
$\epsilon_1-\epsilon_2$ also depends on interaction energies.

The calculation has now supplied a voltage vector for the requested
energy shift. Apply that small step, then locate the transitions
again. Their displacements, converted to energy with the measured
responses, test whether the requested energy changes occurred.
If the chosen $G_{\mathrm P}$ is nearly singular, its inverse
requires large compensating voltages. Another independent pair
of columns of the full $G$ may provide a better set of controls.

\subsection{Changing tunneling at fixed addition energies}

An exchange pulse should strengthen the interdot coupling without
unintentionally moving the dot--reservoir loading thresholds.
We therefore seek a voltage direction that keeps both addition
energies fixed but changes $t_c$. To first order, the first
requirement is simply
\begin{equation}
 G\,d\mathbf V=0.
 \label{eq:calibration-fixed-energy-nullspace}
\end{equation}
For the rank-two matrix above, its nullspace is one-dimensional.
If $G_{\mathrm P}$ is invertible, choose a vector in that
nullspace with barrier component equal to one:
\begin{equation}
 \mathbf n=
 \begin{pmatrix}-G_{\mathrm P}^{-1}\mathbf g_{\mathrm B}\\1\end{pmatrix},
 \qquad
 d\mathbf V=\mathbf n\,dv,\qquad G\mathbf n=0.
 \label{eq:calibration-null-direction}
\end{equation}
Here $dv=dV_{\mathrm B}$ parametrizes the coordinated voltage
change. Equivalently, after choosing the barrier step, the
plunger compensation is
\begin{equation}
 d\mathbf V_{\mathrm P}
 =-G_{\mathrm P}^{-1}\mathbf g_{\mathrm B}\,dV_{\mathrm B}.
 \label{eq:calibration-barrier-compensation}
\end{equation}
The barrier pulse and its two compensating plunger pulses are
applied together. Their energy shifts cancel even though the
potential between the dots, and hence the wave-function overlap,
can change.

For the illustrative matrix above,
\[
 \mathbf n=
 \begin{pmatrix}-17/300\\-13/75\\1\end{pmatrix}.
\]
A $0.30\,\mathrm{mV}$ barrier step therefore requires plunger
steps $-0.017\,\mathrm{mV}$ and $-0.052\,\mathrm{mV}$.
Their combined first-order shifts of both $\mu_1$ and $\mu_2$
vanish. Whether this direction changes tunneling is a separate,
measurable property. Define
$w_j=\partial t_c/\partial V_j$ at the operating point.
Along the coordinated direction,
\begin{equation}
 dt_c=\sum_j w_j\,dV_j
     =\left(\sum_j w_j n_j\right)dv.
 \label{eq:calibration-tunneling-direction}
\end{equation}
A nonzero coefficient permits tunneling control at fixed addition
energies. If it vanishes, these three controls cannot achieve
that change to first order at this operating point.

To measure the coefficient, apply several small values of $v$
in $\mathbf V(v)=\mathbf V^0+\mathbf n v$.
At each setting use the charge-distribution measurement of
Section~\ref{sec:measurement-equilibrium-charge} to determine
$t_c(v)$, and recheck the addition thresholds.
The response must be measured along this coordinated path:
changing the plungers also affects tunneling.

The scan gives $t_c(v)$ along a path that keeps the addition
energies fixed. For spin gates, the same voltage path can be
characterized by its exchange coupling $J(v)$ using
Section~\ref{sec:measurement-exchange-oscillations}.

For several couplings, append response rows such as
$d\ln(t_{ij}/t_{\mathrm{ref}})
=\sum_g B_{ij,g}\,dV_g$, where $t_{\mathrm{ref}}>0$ is a fixed
reference energy. The combined matrix can prescribe a change
of $t_{12}$ with $t_{23}$ and the addition energies held fixed,
provided its rank permits those independent constraints.
Such control of energies and tunnel couplings has been
demonstrated using virtual gates~\cite{Hsiao2020}.

\subsubsection{Turning voltage settings into a pulse}
The energy-shifting and compensated barrier directions solve different
problems: the first loads or unloads electrons, whereas the second
changes their coupling without moving the loading thresholds. A change
between two settings means a change of the physical voltage vector,
not an instantaneous change of a model parameter by itself.
For a locally linear path, a pulse can be programmed as
\[
 \mathbf V(t)=\mathbf V^{\mathrm A}
       +f(t)\bigl(\mathbf V^{\mathrm B}-\mathbf V^{\mathrm A}\bigr),
 \qquad 0\le f(t)\le1.
\]
The function $f$ rises from zero to one, holds, and eventually
returns or proceeds to another setting. For an exchange pulse,
the displacement can follow the compensated direction
$\mathbf n$; for loading, it must instead move an addition energy
across a reservoir chemical potential. The endpoint conditions
depend on whether the operation loads an electron, retains a fixed
charge sector, or converts a spin state into charge motion.

The compensation is local. For a larger excursion, measure or
update the response along the path. The static matrix sets the
voltage combinations; the measured time-dependent response sets
the pulse durations and tests the effect of the edges.


\subsection{Compensating a voltage change across an array}
\label{sec:processor-array-compensation}

Changing the detuning of a target pair may shift the spin gaps
and exchange couplings of neighbouring dots. Near an operating point,
collect the relevant energies and couplings in a vector \(q\),
and the electrode voltages in \(V\). Their measured local
response is
\[
 \delta q_a=\sum_g R_{ag}\delta V_g.
\]
To change one target quantity while holding others fixed,
the requested \(\delta q\) must specify both the target and
the spectator constraints. The requested change must lie in the
range of the measured response matrix. This extends the virtual-gate construction of
Eq.~\eqref{eq:calibration-response-matrix} from two addition energies
to all parameters that the operation must change or preserve.
The response may be short-ranged, but its spatial extent is an
experimental property rather than a consequence of drawing
nearby electrodes.

Consider a square plunger-to-addition-energy response
\(\delta\boldsymbol\mu=G\delta\mathbf V\), with
\(G=D+K\). The diagonal matrix \(D\) describes each
plunger's response on its own dot; \(K\) contains cross
responses. If we use only \(D^{-1}\), the voltage
\(D^{-1}\delta\boldsymbol\mu\) produces an extra energy
shift \(KD^{-1}\delta\boldsymbol\mu\). For
\(\|D^{-1}K\|<1\), a Neumann expansion cancels these
successive cross effects:
\begin{equation}
 G^{-1}
 =[I-D^{-1}K+(D^{-1}K)^2-\cdots]D^{-1}.
 \label{eq:processor-local-compensation}
\end{equation}
The equivalent iteration is
\[
 \delta\mathbf V^{(k+1)}
 =D^{-1}\bigl(\delta\boldsymbol\mu-K\delta\mathbf V^{(k)}\bigr).
\]
If \(K\) connects only neighbouring dots, each additional
order reaches farther through the array. When the response is
strong or long-ranged, the truncated local inverse may be
inadequate and the wider matrix must be used.

This inverse is only as current as its measured response.
Changing a barrier can alter the two neighbouring spin gaps,
their exchange, and the reference distribution of a shared
sensor. Before reusing a gate pulse, recheck the observables
that depend on that barrier; a changed sensor distribution
must be separated from a changed gate outcome. A feedback
procedure for updating calibrated controls is developed in
Eq.~\eqref{eq:calibration-feedback-update}. Its use in an array
also requires tracking which other calibrated operations depend
on the control being updated.
